\documentclass[11pt]{article}

\usepackage[margin=1in]{geometry}
\usepackage{amsmath,amssymb,amsfonts,bm,mathtools}
\usepackage{booktabs}
\usepackage{enumitem}
\usepackage[backend=biber,style=numeric,sorting=none,natbib=true]{biblatex}
\usepackage{hyperref}
\usepackage{xcolor}
\usepackage{graphicx}
\usepackage{multirow}
\usepackage{algorithm}
\usepackage{algpseudocode}

\hypersetup{
    colorlinks=true,
    linkcolor=blue,
    citecolor=blue,
    urlcolor=blue
}

\addbibresource{references.bib}

\title{Certifiable Calibration Data Design via Information-Theoretic Subset Selection}
\author{}
\date{}

\begin{document}

\maketitle

\begin{abstract}
Calibration quality in robotics depends not only on the estimation backend, but
also on the informativeness of the data used for calibration. In practice,
calibration datasets are still commonly collected using heuristics such as
coverage, view diversity, or operator intuition. While these heuristics are
often useful, they do not quantify whether a selected dataset sufficiently
constrains the calibration parameters of interest. This paper formulates
calibration data collection as a finite-budget experimental-design problem over
candidate views, poses, or trajectory segments. We use the marginal Fisher
information matrix of the calibration parameters as the design object and study
E-, A-, and D-optimal subset selection. To move beyond heuristic selection, we
combine greedy and Frank--Wolfe optimization with post-hoc certification from
the relaxed objective value, which upper-bounds the best achievable discrete
objective under the same linearization model. Experiments in Isaac Sim, on two
real ChArUco capture sessions, and with two downstream backends show a nuanced
picture. Information-aware subset selection consistently outperforms random and
geometric coverage baselines. Exhaustive search on reduced candidate pools shows
that the relaxation provides a meaningful certificate on discrete subset
quality. At the same time, the best design criterion depends on the downstream
solver: E-optimal selection is effective when the linearized information model
tracks estimator behavior, whereas A-optimal selection is more robust for
bundle-adjustment backends such as Kalibr. These results suggest that
calibration should be treated as a data-design problem, but that certification
must be interpreted relative to the adopted information model and solver
assumptions.
\end{abstract}

\section{Introduction}
Robotic perception systems depend critically on accurate calibration.
Errors in camera intrinsics, sensor extrinsics, temporal offsets, or related
latent parameters propagate directly into visual odometry, SLAM, sensor fusion,
and control. For that reason, calibration has traditionally been treated as a
central estimation problem in robotics and computer vision.

In practice, however, calibration success depends on more than the optimizer.
Even strong calibration backends fail or become unstable when the collected data
weakly excites important parameter directions. Poor parallax, repeated
near-identical views, target placements concentrated near the image center, or
trajectory segments that leave certain modes nearly unobservable can all yield
fragile estimates despite low training reprojection error. Practitioners
therefore face a second problem that is often treated informally: how to collect
the data that makes calibration well-conditioned in the first place.

Most commonly used pipelines do not explicitly solve this upstream design
problem. Classical calibration methods such as Zhang's planar approach and
practical robotics toolchains such as Kalibr are highly effective once a
dataset has already been collected
\citep{zhang2000flexible,furgale2013kalibr,kalibr_toolbox}. Active and
next-best-view methods improve the collection process by proposing new poses
based on uncertainty reduction, visibility, or expected information gain
\citep{richardson2013aprilcal,peng2019calibration,rojtberg2018efficient,yang2023nbv,wang2025active}.
Observability-aware self-calibration work has likewise shown that not all motion
segments are equally useful
\citep{maye2013self,schneider2017informative,schneider2019observability,yang2023online}.
Yet a gap remains between ``better data'' and ``certifiably informative data'':
most prior methods are sequential, heuristic, or solver-specific, rather than a
global subset-selection formulation with a quantitative certificate on the final
selected set.

This paper takes a step toward that goal. We formulate calibration data
collection as a finite-budget subset-selection problem over a candidate pool of
views, poses, or motion segments. Using the marginal Fisher information matrix
of the calibration parameters, we optimize standard experimental-design
criteria, with particular emphasis on E-optimality, which strengthens the
weakest observable direction. We then use a continuous relaxation to obtain a
computable upper bound on the best achievable discrete objective under the same
linearization model. This does \emph{not} certify global optimality of the
nonlinear calibration estimate itself, nor does it imply that every downstream
backend will prefer the same subset. What it does provide is a clear,
model-relative statement about how good the selected data is under a stated
design surrogate.

The main message of the paper is therefore more precise than ``E-optimality is
always best.'' Instead, we argue that:
\begin{enumerate}[leftmargin=1.5em]
    \item calibration should be treated as a data-design problem in addition to
    an estimation problem;
    \item information-theoretic subset selection provides a principled way to
    rank candidate calibration data;
    \item continuous relaxations provide meaningful post-hoc certificates on the
    adopted design objective; and
    \item the best design criterion depends on the mismatch between the
    linearized information surrogate and the downstream solver.
\end{enumerate}

This framing is especially relevant for robotics, where calibration often has to
be performed under limited operator time, motion budget, or sensing constraints.
A method that can say not just ``this dataset seems reasonable,'' but ``this
dataset is near-optimal under a stated design model'' is useful even when the
certificate is explicitly model-relative rather than absolute.

\section{Related Work}
\subsection{Calibration Backends}
Classical calibration methods estimate intrinsics and extrinsics from a fixed
dataset. Zhang's planar-pattern method remains a foundational reference for
practical camera calibration \citep{zhang2000flexible}, while Kalibr has become
a widely used robotics backend for multi-camera and visual-inertial calibration
\citep{furgale2013kalibr,kalibr_toolbox}. These are estimation backends, not
data-design methods.

\subsection{Observability and Informative Data}
Observability-aware calibration work has shown that calibration quality depends
strongly on excitation and motion diversity. Information-theoretic
self-calibration, informative segment retention, and observability analysis for
visual-inertial systems all support the view that weak datasets are a primary
source of failure
\citep{maye2013self,schneider2017informative,schneider2019observability,yang2023online}.
Our formulation is aligned with this literature, but focuses on fixed-budget
subset selection with post-hoc certification.

\subsection{Active Calibration and Guided Capture}
AprilCal, Calibration Wizard, and next-best-view systems explicitly guide the
collection process
\citep{richardson2013aprilcal,peng2019calibration,rojtberg2018efficient,yang2023nbv,wang2025active}.
These methods are highly relevant because they recognize the importance of the
data-acquisition stage. The main difference in our setting is that we optimize
over a finite candidate pool and ask for a certificate on the quality of the
final selected subset rather than a sequential policy alone.

\subsection{Certifiable Robotics and OASIS}
Recent robotics work has used relaxation-based tools to provide certificates for
either estimation or design. Certifiable calibration solvers focus on globally
solving a fixed calibration problem \citep{giamou2019certifiable}, whereas
OASIS formulates sensor placement for SLAM as an E-optimal subset-selection
problem with verification through convex relaxation \citep{kaveti2023oasis}. Our
work is conceptually closest to the latter: the goal is not to certify the
nonlinear calibration estimate itself, but to certify the informativeness of the
selected data under a stated surrogate model.

\subsection{Learning-Based Calibration}
Learning-based approaches such as RegNet, CalibNet, and LCCNet improve
automation and targetless operation
\citep{schneider2017regnet,iyer2018calibnet,lv2021lccnet}. Recent surveys and
critiques, however, continue to report challenges in generalization and
evaluation consistency \citep{liao2023survey,huang2025whatmatters}. These
methods address a different question from ours: they improve how calibration is
predicted or refined, but do not directly provide observability certificates on
the collected dataset.

\section{Problem Formulation}
Let
\begin{equation}
\mathcal{C} = \{c_1,\dots,c_N\}
\end{equation}
be a finite set of candidate calibration data units. A candidate may correspond
to a single image, a target pose, a robot configuration, or a short trajectory
segment. We seek a subset of size $K$ that is maximally informative for a set
of calibration parameters $\theta \in \mathbb{R}^p$.

Let $\eta$ denote nuisance variables such as per-frame target pose or camera
pose. For candidate $c_i$, let $r_i(\theta,\eta)$ be the residual vector and
let $\Sigma_i$ be the measurement covariance. Under a Gauss--Newton
approximation, the candidate-level Fisher information is
\begin{equation}
\mathcal{I}_i
\approx
J_i^\top \Sigma_i^{-1} J_i,
\qquad
J_i = \begin{bmatrix} J_{i,\theta} & J_{i,\eta} \end{bmatrix}.
\end{equation}
Because the information contributions add across conditionally independent
measurements,
\begin{equation}
\mathcal{I}(u) = \sum_{i=1}^{N} u_i \mathcal{I}_i,
\qquad
u \in \{0,1\}^N,
\qquad
\mathbf{1}^\top u = K.
\end{equation}

Partitioning the full information matrix into calibration and nuisance blocks,
\begin{equation}
\mathcal{I}(u) =
\begin{bmatrix}
\mathcal{I}_{\theta\theta}(u) & \mathcal{I}_{\theta\eta}(u) \\
\mathcal{I}_{\eta\theta}(u) & \mathcal{I}_{\eta\eta}(u)
\end{bmatrix},
\end{equation}
we score only the calibration parameters through the Schur complement
\begin{equation}
\widetilde{\mathcal{I}}_\theta(u)
=
\mathcal{I}_{\theta\theta}(u)
- \mathcal{I}_{\theta\eta}(u)\mathcal{I}_{\eta\eta}(u)^{-1}\mathcal{I}_{\eta\theta}(u).
\label{eq:schur}
\end{equation}
This marginal information matrix is the design object used throughout the paper.

We study three standard criteria:
\begin{align}
\text{E-optimal: } &
\max_{u}
\lambda_{\min}\!\bigl(\widetilde{\mathcal{I}}_\theta(u)\bigr),
\\
\text{A-optimal: } &
\min_{u}
\operatorname{tr}\!\bigl(\widetilde{\mathcal{I}}_\theta(u)^{-1}\bigr),
\\
\text{D-optimal: } &
\max_{u}
\log\det\!\bigl(\widetilde{\mathcal{I}}_\theta(u)\bigr),
\end{align}
all under the binary budget constraint above. E-optimality is attractive
because it explicitly strengthens the weakest observable direction, but one of
the conclusions of this paper is that no single criterion dominates across all
solvers and noise regimes.

\section{Method}
\subsection{Greedy Discrete Selection}
The simplest baseline is greedy subset construction. Starting from the empty
set, each step adds the candidate that most improves the chosen design
criterion. This yields practical discrete baselines for E-, A-, and D-optimal
selection.

\subsection{Continuous Relaxation}
To obtain a certificate, we relax the binary selection variables to continuous
weights:
\begin{equation}
u \in [0,1]^N,
\qquad
\mathbf{1}^\top u = K.
\label{eq:relax}
\end{equation}
This feasible set strictly contains the binary feasible set, so the relaxed
optimum upper-bounds the best discrete objective for maximization problems and
lower-bounds it for minimization problems. For E-optimality, the relaxed problem
is
\begin{equation}
\max_{u \in [0,1]^N,\ \mathbf{1}^\top u = K}
\lambda_{\min}\!\bigl(\widetilde{\mathcal{I}}_\theta(u)\bigr).
\label{eq:e_relax}
\end{equation}
We optimize the relaxation with Frank--Wolfe.

\subsection{Rounding and Certification}
After optimization, the continuous solution is rounded to a binary subset by
selecting the $K$ largest weights. Let $f_{\mathrm{relax}}$ be the relaxed
objective and let $f_{\mathrm{disc}}$ be the objective of the rounded subset.
For maximization criteria, the relaxed value is an upper bound on the unknown
integer optimum $f^\star$:
\begin{equation}
f_{\mathrm{disc}} \leq f^\star \leq f_{\mathrm{relax}}.
\end{equation}
This yields a certificate gap
\begin{equation}
\Delta_{\mathrm{cert}} = f_{\mathrm{relax}} - f_{\mathrm{disc}},
\end{equation}
or a normalized variant if desired. A small gap means the rounded subset is
near-optimal under the chosen information model.

This point is crucial: the certificate applies to the \emph{design surrogate},
not to nonlinear calibration accuracy in general. It does not certify global
optimality of the downstream estimator, nor does it guarantee that all
backends will prefer the same subset.

\subsection{Practical Computation}
In the experiments, per-candidate information blocks are computed at a fixed
bootstrap linearization point. The nuisance block is block-diagonal because each
candidate contributes an independent pose block, so the Schur complement can be
formed efficiently. The Frank--Wolfe linear minimization oracle under the
budget-constrained box relaxation has a closed-form top-$K$ solution, which
replaces a generic LP solve. In the final experiments we use a conservative
stopping rule requiring both small relative objective change and small duality
gap; this avoids premature stopping but does not eliminate small-budget
rounding issues.

\begin{algorithm}[ht]
\caption{Certifiable Calibration Data Selection}
\begin{algorithmic}[1]
\Require Candidate pool $\mathcal{C}$, budget $K$, linearization point
$(\bar{\theta},\bar{\eta})$, design criterion $J$
\State Compute per-candidate information blocks $\mathcal{I}_i$
\State Build $\widetilde{\mathcal{I}}_\theta(u)$ via the Schur complement
\State Optimize the relaxed problem for criterion $J$ with Frank--Wolfe
\State Round the final continuous weights to a binary top-$K$ subset
\State Evaluate the discrete objective and certificate gap
\State Return selected subset, objective value, and certificate gap
\end{algorithmic}
\end{algorithm}

\subsection{Scope}
The method makes three assumptions that must be explicit.
First, informativeness is measured relative to a chosen linearization point.
Second, certification applies to the information-theoretic surrogate, not to
arbitrary nonlinear solver behavior. Third, a dataset that is strong for one
backend need not be optimal for another. We therefore treat the certificate as
model-relative evidence rather than an absolute claim about all calibration
pipelines.

\section{Experimental Protocol}
\subsection{Questions}
We evaluate five questions:
\begin{enumerate}[leftmargin=1.5em]
    \item Do information-theoretic subsets outperform random and coverage
    baselines?
    \item Does the relaxed objective provide a meaningful certificate on subset
    quality?
    \item How sensitive are conclusions to noise, distortion, and subset budget?
    \item Do the conclusions persist in real data?
    \item How strongly do the conclusions depend on the downstream calibration
    backend?
\end{enumerate}

\subsection{Simulation Setup}
The main synthetic benchmark uses $N=720$ candidate poses generated in Isaac
Sim on a hemisphere around a fixed target. The default budget is $K=20$. Two
targets are compared under matched geometry: an AprilGrid and a ChArUco board
with identical corner count, pitch, and physical extent. Noise is added as
i.i.d.\ Gaussian image perturbation at near-noiseless, moderate, and high
levels. Additional experiments introduce medium and high Brown--Conrady
distortion. The main downstream backend is OpenCV
\texttt{cv2.calibrateCamera}, which lets us compare all selection methods under
a common solver.

\subsection{Methods and Metrics}
We compare random subset selection, geometric coverage via farthest-point
sampling in translation space, greedy E-optimal, A-optimal, and D-optimal
selection, and Frank--Wolfe E-optimal, A-optimal, and D-optimal selection.

We report both design and calibration metrics. Design quality is measured using
$\lambda_{\min}(\widetilde{\mathcal{I}}_\theta)$ and $\log\det
\widetilde{\mathcal{I}}_\theta$. Calibration quality is measured using focal
length error, principal-point error, distortion-parameter error, training
reprojection RMS, and held-out reprojection RMS. We intentionally avoid a
single aggregate intrinsic-parameter norm because it mixes pixel-valued and
dimensionless quantities in a way that is hard to interpret physically.

\subsection{Real-World Validation}
We use two real ChArUco sessions. The first is a handheld capture with
$N=76$ retained frames and a 20-frame held-out split. The second is a more
challenging session with stronger out-of-plane tilts and multiple depths,
yielding $N=66$ candidate frames after filtering. In the second session,
held-out poses are fixed once using all-frame intrinsics, specifically to avoid
per-method re-fitting that could otherwise mask intrinsic differences.

\subsection{Alternative Backend}
To test solver dependence, we repeat the evaluation with Kalibr on the noiseless
Isaac Sim AprilGrid benchmark. Here each method selects from the same 700-view
candidate pool after reserving 20 held-out views. Kalibr performs joint
bundle-adjustment over intrinsics and per-view extrinsics, providing a strong
contrast to the OpenCV fixed-backend experiments.

\section{Results}
\subsection{Information-Aware Selection Beats Heuristics}
Across both targets and all tested noise levels, information-aware selection
consistently outperforms random subsets and geometric coverage. This gap is
especially visible in held-out error at moderate and high noise, where coverage
can look geometrically diverse while still leaving important calibration
directions weakly constrained.

One particularly robust finding is that training reprojection RMS varies little
across methods even when held-out performance differs substantially. In other
words, low training residual is not a reliable indicator that the chosen dataset
is informative. This is an important practical message independent of which
design criterion ultimately wins.

\subsection{No Single Criterion Dominates}
The OpenCV-based experiments show a consistent but nuanced pattern.
At moderate noise on AprilGrid, FW-E attains the highest
$\lambda_{\min}$ ($0.874$), Greedy-E achieves the best focal error
($0.150$ px), and FW-A achieves the best held-out RMS ($0.819$ px).
At high noise, A-optimal becomes the most reliable criterion for held-out
accuracy, with the best held-out RMS on both AprilGrid ($1.782$ px) and matched
ChArUco ($2.483$ px).

This is the right place to narrow the paper's claim. The evidence does not
support ``E-optimality universally wins.'' What it supports is that
information-theoretic design is useful, that E-optimality is valuable when the
weakest mode is the main concern, and that A-optimality is often more robust
when noise or solver coupling makes balanced conditioning more important than
the single weakest direction.

\subsection{The Certificate Is Meaningful for the Design Objective}
Exhaustive search on reduced candidate pools shows that the relaxed objective is
a useful certificate on discrete subset quality. On the tested small-$N$
instances, Frank--Wolfe E-optimal rounding comes within $0$--$2.5\%$ of the
true discrete optimum, whereas greedy E-optimal selection is typically
$7$--$19\%$ below optimum. The relaxed value remains a valid upper bound on the
unknown integer optimum under the same surrogate.

These results do not prove that the same near-optimality gap holds at large
scale, but they do support the claim that the relaxation is more than a
numerical convenience. It yields a real, interpretable certificate on the
subset-selection problem being solved.

\subsection{Budget and Distortion Effects}
The budget sweep shows strong diminishing returns beyond roughly
$K \approx 20$--$30$ in the tested settings. This is useful operationally:
substantial gains are realized with relatively small datasets, but there is
little benefit in continuing to add near-redundant views.

At the same time, the budget sweep exposes an important limitation. Continuous
E-optimal relaxation can round poorly at small budgets because the solution
remains too diffuse. This is why A-optimal selection is the safer default when
$K$ is very small, even if E-optimality becomes competitive or superior once
the budget is large enough for the relaxation to concentrate.

The distortion experiments are encouraging. Under both medium and high
Brown--Conrady distortion, the qualitative ranking among methods remains similar
to the zero-distortion case. This suggests that the subset-selection behavior is
not a fragile artifact of the ideal pinhole setting alone.

\subsection{Real-World Validation Supports the Main Trend}
The two real ChArUco sessions are smaller than the synthetic benchmark, but
they are still informative.

In the first session, Greedy-E achieves the best held-out RMS ($0.752$ px),
followed closely by FW-E ($0.764$ px), while coverage is worst ($1.044$ px).
In the second, more honest fixed-pose evaluation, A-optimal and FW-A achieve
the best held-out RMS ($0.737$ px), Greedy-E follows at $0.772$ px, and random
selection is much worse at $1.963$ px. The most important real-world conclusion
is therefore not which FIM criterion wins by a small margin, but that
information-aware selection consistently beats random and geometric heuristics.

\begin{table}[ht]
\centering
\small
\caption{Representative real-world results.
Session 1 uses a standard held-out reprojection protocol.
Session 3 fixes held-out poses once across all methods to avoid per-method
re-fitting effects.}
\label{tab:real_world_summary}
\setlength{\tabcolsep}{5pt}
\begin{tabular}{lrrr}
\toprule
Setting & Best method(s) & Best held-out RMS & Baseline comparison \\
\midrule
Session 1 & Greedy-E & $0.752$ px & Coverage: $1.044$ px \\
Session 3 & A-optimal, FW-A & $0.737$ px & Random: $1.963$ px \\
\bottomrule
\end{tabular}
\end{table}

\subsection{Backend Sensitivity Is the Paper's Sharpest Result}
The strongest cautionary result comes from the Kalibr backend. At
$K=20$, A-optimal achieves a held-out RMS of $0.090$ px, whereas Greedy-E and
FW-E degrade to $0.732$ px and $0.761$ px respectively. This is not a small
reordering. It is a qualitative reversal of the OpenCV-style narrative.

Inspection of the selected subsets explains why. Under the adopted surrogate,
E-optimal selection can favor a tightly clustered set of views that strongly
constrain the weakest mode of the marginal information matrix. Kalibr, however,
jointly refines intrinsics and per-view extrinsics. It benefits from broader
angular spread across the hemisphere to keep the pose subspace well-conditioned.
That behavior is captured more naturally by A-optimality than by a pure
weakest-mode objective.

This result should not be hidden because it actually clarifies the contribution.
The certificate is meaningful relative to the adopted surrogate model; it is
not a guarantee that every downstream solver will prefer the same subset. For
backends whose behavior departs substantially from the local information model,
A-optimality or a geometry-aware hybrid criterion may be the better practical
choice.

\subsection{A Practical Criterion Guide}
The experiments suggest three diagnostic axes for choosing a criterion:
solver type, noise regime, and budget.

If the downstream solver conditions on known or pre-solved poses, E-optimality
is attractive once the budget is large enough for reliable rounding. If the
solver jointly optimizes intrinsics and per-view extrinsics, A-optimality is the
safer choice because it promotes broader global conditioning. Likewise,
A-optimality is safer at high noise and at small budgets, whereas E-optimality
is most compelling for fixed-pose solvers at moderate noise and
$K \ge 20$.

\begin{table}[ht]
\centering
\small
\caption{Practical criterion selection guide.}
\label{tab:criterion_guide}
\setlength{\tabcolsep}{5pt}
\begin{tabular}{llll}
\toprule
Solver type & Noise level & Budget & Recommended criterion \\
\midrule
Fixed-pose solver & Low to moderate & $K < 15$ & A-optimal \\
Fixed-pose solver & Low to moderate & $K \ge 20$ & FW-E or Greedy-E \\
Fixed-pose solver & High noise & Any & A-optimal \\
Joint BA solver & Any & Any & A-optimal \\
\bottomrule
\end{tabular}
\end{table}

\section{Discussion}
Three conclusions stand out.

First, calibration data selection should be treated as an optimization problem
rather than an informal capture step. Even simple information-theoretic subset
selection beats strong heuristic baselines in many settings.

Second, certification is possible, but only in a model-relative sense. The
continuous relaxation provides a valid bound for the adopted design problem, not
for arbitrary nonlinear calibration behavior. This still has practical value:
it lets a user compare subsets with a stated confidence about near-optimality
under the chosen surrogate.

Third, solver mismatch matters. This is not a side detail. It is the main
boundary of the present approach. If the downstream backend behaves similarly to
the linearized information model, E-optimality is attractive because it
strengthens the weakest mode directly. If the backend relies on broader
geometric diversity, a criterion that balances information globally, such as
A-optimality, can be more robust.

The current evidence also defines clear limitations.
The certification claim is relative to the design surrogate, not the nonlinear
estimator. Most experiments are still synthetic. The real-world validation is
target-based and monocular rather than a broad cross-sensor study. Small-budget
rounding remains a known weakness of the current continuous E-optimal pipeline.
These are not peripheral caveats; they define the present scope of what is
supported.

\section{Conclusion}
This paper reformulated calibration data collection as a certifiable
subset-selection problem over candidate measurements. By combining Fisher
information-based design criteria with continuous relaxation and post-hoc
certification, it provided a principled alternative to heuristic calibration
capture. The experiments show that information-aware selection improves over
random and coverage-based baselines, that the relaxed objective provides a
meaningful certificate on the adopted design problem, and that the best
criterion depends strongly on the downstream backend. The broader lesson is
that certification in calibration should be stated carefully: it is strongest
when the design surrogate, solver, and evaluation protocol are aligned. Future
work should focus on backend-aware design criteria, stronger real-world
validation, and tighter rounding schemes for low-budget regimes.

\paragraph{Submission note.}
This revised draft intentionally trims the table-heavy presentation of the
working manuscript and sharpens the claims to match the evidence. The full
K-sweep, distortion, and per-metric tables are better treated as supplementary
material in a conference submission.

\printbibliography

\end{document}
